\notechapter{N-20220811}{Inequality Exercises: Cauchy, AM--GM, Vectors, and Optimization}

\section{How to read this note}

This note is a compact inequality practice sheet.  The solutions are handwritten, but the recurring pattern is clear: almost every problem is asking us to recognize which standard inequality is hiding behind the expression.  The useful question is not ``which formula should I memorize?'', but rather ``what structure does the expression already have?''

If there are fractions with variable denominators, Cauchy or Engel form is natural.  If there are products such as \(abc\), AM--GM is natural.  If there are square roots of quadratic expressions, think of vector lengths and the triangle inequality.  If the problem asks for a maximum or minimum under a quadratic constraint, expect equality cases in Cauchy--Schwarz.

\section{A fraction problem with a constraint}

One problem asks to handle expressions of the form
\[
  \frac{4}{4-x^2}+\frac{9}{9-y^2}
\]
under constraints such as \(xy=-1\) and bounded intervals for \(x,y\).  The handwritten solution uses a Cauchy-style lower bound:
\[
  \frac{4}{4-x^2}+\frac{9}{9-y^2}
  \ge
  \frac{(2+3)^2}{(4-x^2)+(9-y^2)}.
\]
Then the constraint controls the denominator.  The important point is that the numerator \(4,9\) are squares, so the expression is inviting the Cauchy inequality
\[
  \frac{a^2}{u}+\frac{b^2}{v}\ge \frac{(a+b)^2}{u+v}.
\]

This is one of the most common contest moves: when a sum of fractions has square numerators, try to combine them by Cauchy rather than simplifying separately.

\section{A weighted square-root comparison}

Another exercise defines
\[
  P=\sqrt{ab}+\sqrt{cd},
\]
and
\[
  Q=\sqrt{ma+nc}\sqrt{\frac bm+\frac dn},
\]
where the variables are positive.  Expanding \(Q^2\) gives
\[
\begin{aligned}
Q^2
&=(ma+nc)\left(\frac bm+\frac dn\right)\\
&=ab+cd+\frac{mad}{n}+\frac{ncb}{m}.
\end{aligned}
\]
By AM--GM,
\[
  \frac{mad}{n}+\frac{ncb}{m}\ge 2\sqrt{abcd}.
\]
Hence
\[
  Q^2\ge ab+cd+2\sqrt{abcd}=(\sqrt{ab}+\sqrt{cd})^2=P^2.
\]
So \(P\le Q\), with equality exactly when
\[
  \frac{mad}{n}=\frac{ncb}{m}.
\]
The original note asks for the condition for strict inequality; it is precisely the negation of this equality condition.

\section{Maximizing a radical expression}

The note then asks for the maximum of
\[
  f(x)=2\sqrt{x-3}+\sqrt{5-x},\qquad 3\le x\le5.
\]
One way is calculus.  We have
\[
  f'(x)=\frac1{\sqrt{x-3}}-\frac1{2\sqrt{5-x}}.
\]
Setting \(f'(x)=0\) gives
\[
  2\sqrt{5-x}=\sqrt{x-3},
\]
so
\[
  4(5-x)=x-3,\qquad x=\frac{23}{5}.
\]
At this point,
\[
  f\left(\frac{23}{5}\right)
  =2\sqrt{\frac85}+\sqrt{\frac25}
  =\sqrt{10}.
\]

But the handwritten note also asks whether this can be done without derivatives.  The answer is yes, by Cauchy--Schwarz:
\[
  \left(2\sqrt{x-3}+\sqrt{5-x}\right)^2
  \le (4+1)((x-3)+(5-x))=10.
\]
Equality holds when
\[
  \frac{2}{1}=\frac{\sqrt{x-3}}{\sqrt{5-x}},
\]
which again gives \(x=23/5\).  This non-calculus proof is more elegant because it explains the form of the answer.

\section{Equality in Cauchy--Schwarz}

A vector problem gives
\[
  a^2+b^2+c^2=25,\qquad x^2+y^2+z^2=36,\qquad ax+by+cz=30.
\]
By Cauchy--Schwarz,
\[
  (ax+by+cz)^2\le(a^2+b^2+c^2)(x^2+y^2+z^2)=25\cdot36=900.
\]
Since \(ax+by+cz=30\), equality holds.  Therefore
\[
  (a,b,c)=k(x,y,z).
\]
Comparing norms gives
\[
  k=\frac56.
\]
Thus
\[
  \frac{a+b+c}{x+y+z}=\frac56.
\]
This is an excellent example of why equality cases matter.  The inequality is not just bounding a number; it tells us the exact geometry of the vectors.

\section{Triangle-angle inequalities}

Let \(A,B,C\) be the angles of a triangle.  Then
\[
  A+B+C=\pi.
\]
By Cauchy--Schwarz,
\[
  (A+B+C)\left(\frac1A+\frac1B+\frac1C\right)\ge(1+1+1)^2=9.
\]
Hence
\[
  \frac1A+\frac1B+\frac1C\ge\frac9\pi.
\]
Similarly,
\[
  A^2+B^2+C^2\ge\frac{(A+B+C)^2}{3}=\frac{\pi^2}{3}.
\]
Both equalities occur only for
\[
  A=B=C=\frac\pi3.
\]
The note's implicit lesson is that symmetric triangle inequalities often have equality at the equilateral triangle.

\section{Two spheres and a linear equation}

Another problem asks to solve a system of two sphere equations.  The trick is to subtract them.  If two equations are
\[
  (x-a_1)^2+(y-b_1)^2+(z-c_1)^2=R_1^2,
\]
\[
  (x-a_2)^2+(y-b_2)^2+(z-c_2)^2=R_2^2,
\]
subtracting cancels \(x^2+y^2+z^2\) and leaves a plane.  The intersection of two spheres is therefore a circle, a point, or empty; algebraically it becomes a sphere-plane problem.

The handwritten solution then uses Cauchy--Schwarz to bound a linear expression such as
\[
  -8x+6y-24z.
\]
The equality condition determines the direction of the vector \((x,y,z)\).  Again the method is geometric: a linear function is maximized on a sphere in the direction of its coefficient vector.

\section{A product inequality}

If
\[
  (1+a)(1+b)(1+c)=8,\qquad a,b,c>0,
\]
then the note proves
\[
  abc\le1.
\]
Indeed,
\[
\begin{aligned}
8&=1+a+b+c+ab+bc+ca+abc\\
&\ge 1+3\sqrt[3]{abc}+3\sqrt[3]{(ab)(bc)(ca)}+abc\\
&=1+3t+3t^2+t^3=(1+t)^3,
\end{aligned}
\]
where \(t=\sqrt[3]{abc}\).  Thus \(1+t\le2\), so \(t\le1\), and therefore \(abc\le1\).

\section{Vector-style radical inequality}

The pages also contain an inequality of the form
\[
  \sqrt{x^2+xy+y^2}+\sqrt{x^2+xz+z^2}\ge \sqrt{y^2+yz+z^2}.
\]
This is best understood as a disguised triangle inequality.  One rewrites each quadratic as a squared Euclidean norm of a two-dimensional vector.  For example,
\[
  x^2+xy+y^2=\left(x+\frac y2\right)^2+\frac34y^2.
\]
Once each radical is interpreted as a vector length, the inequality becomes
\[
  \|u\|+\|v\|\ge\|u+v\|.
\]
This is a typical high-level move: turn algebraic square roots into geometry.

\section{Expanded synthesis and advanced context}

The note is not merely a list of inequality tricks.  It repeatedly uses the same three ideas: Cauchy detects equality of vectors, AM--GM detects equality of positive quantities, and radical inequalities often hide triangle inequalities.  The more mature way to remember the sheet is therefore:
\[
  \text{fractions}\to\text{Cauchy},\qquad
  \text{products}\to\text{AM--GM},\qquad
  \text{radicals}\to\text{norms}.
\]



\paragraph{Expanded synthesis.}
For this chapter, the elementary material should be read as a study of what remains invariant when coordinates change. The earlier sections (`How to read this note`, `A fraction problem with a constraint`, `A weighted square-root comparison`, `Maximizing a radical expression`, `Equality in Cauchy--Schwarz`) compute with arrays, bases, pivots, determinants, or eigenvectors; the deeper object is a linear map together with the spaces it acts on. A matrix is a representation of that map after bases have been chosen. This is why formulas such as
\[
  A' = P^{-1}AP,\qquad \widetilde A = T^{-1}AS
\]
are not cosmetic: they distinguish an intrinsic morphism from a coordinate description.

The structural hierarchy is as follows. Kernels record what a map forgets, images record what it reaches, cokernels record obstructions to solvability, and quotients record information after an equivalence has been imposed. Determinants act on the top exterior power and therefore measure oriented volume; traces are invariant under cyclic rearrangement and become characters in representation theory; adjoints depend on an inner product and lead to self-adjoint spectral theory.

A useful way to return to the computations is to ask, for every formula, which invariant it is measuring. Row reduction measures rank and kernel; diagonalization measures decomposition into invariant subspaces; Gram--Schmidt measures orthogonal projection; positive definiteness measures ellipsoid geometry; tensor notation measures how components transform. The elementary methods are therefore the finite-dimensional laboratory in which duality, quotienting, functoriality, and spectral decomposition can be seen clearly.


\section{Elementary-to-advanced bridge: equality cases as geometry}

The original inequality exercises repeatedly ask when equality holds.  This is not a side detail; equality describes the geometry of the inequality.  In Cauchy--Schwarz,
\[
  \left(\sum a_ib_i\right)^2\le \left(\sum a_i^2\right)\left(\sum b_i^2\right),
\]
equality holds exactly when the vectors $(a_i)$ and $(b_i)$ are linearly dependent.  In AM--GM,
\[
  \frac{x_1+\cdots+x_n}{n}\ge \sqrt[n]{x_1\cdots x_n},
\]
equality holds when all variables are equal.  These equality sets often reveal the hidden symmetry of a problem.

\section{Elementary-to-advanced bridge: Lagrange multipliers behind constrained inequalities}

Many inequality exercises are optimization problems with constraints.  If one wants to extremize $f(x_1,\ldots,x_n)$ under $g(x_1,\ldots,x_n)=0$, the smooth interior candidates satisfy
\[
  \nabla f=\lambda \nabla g.
\]
For example, optimizing a quadratic form under $\|x\|=1$ leads to
\[
  Ax=\lambda x,
\]
so eigenvalues are extrema of quadratic forms on the unit sphere.  This is the Rayleigh quotient principle:
\[
  \lambda_{\min}\le \frac{x^TAx}{x^Tx}\le \lambda_{\max}
\]
for symmetric $A$.

\section{Elementary-to-advanced bridge: semidefinite viewpoint}

Some inequalities are equivalent to positive semidefiniteness.  To prove
\[
  ax^2+2bxy+cy^2\ge0
\]
for all $x,y$, one checks
\[
  \begin{pmatrix}a&b\\b&c\end{pmatrix}\succeq0.
\]
That means
\[
  a\ge0,\qquad ac-b^2\ge0.
\]
This connects elementary quadratic inequalities to semidefinite programming, where one optimizes linear functions over the cone of positive semidefinite matrices.

\section{Returning to the original problems}

The contest techniques in the source—Cauchy, AM--GM, substitutions, and equality checking—should be read as finite-dimensional optimization.  The advanced language does not replace the computations; it explains why the same few patterns recur: projection, convexity, eigenvalues, and positive semidefinite forms.
